% From A3_Energetic0.tex, with revisions
% for ixM
\subsection{Shear (Energetic) Trace}\label{sec:trace-energy}

The energy-conjugate trace is defined by requiring that the beam strains 
$\bm{e} = (\epsilon, \bm{\gamma}_\perp, \varkappa, \bm{\kappa}_\perp)$ 
be work-conjugate to the classical resultants 
$\bm{s} = (N, \bm{F}, T, \bm{M})$
\eqref{eq:def-trace-conjugate}:
\begin{equation*}
\delta\bm{e}\cdot \bm{s} = \int_\Section \delta\bm{\varepsilon} : \bm{\sigma} \dA
\end{equation*}
for all variations within the Saint~Venant class.

We seek beam strains of the form $\bm{e}(x) = \mathbf{X}\hat{\bm{e}}(x)$ with 
$\mathbf{X}$ independent of $x$. 
The energy conjugacy condition 
\eqref{eq:def-trace-conjugate} requires:
\begin{equation}
\bm{e}\cdot \bm{s} = \hat{\bm{e}}\cdot \bar{\mathbf{C}} \hat{\bm{e}}.
\end{equation}

Substituting $\bm{e} = \mathbf{X}\hat{\bm{e}}$ and $\bm{s} = -\mathbf{R}\mathbf{L}^{-1}\hat{\bm{e}}$:
\begin{equation*}
-\hat{\bm{e}}\cdot \mathbf{X}^\top \mathbf{R} \mathbf{L}^{-1} \hat{\bm{e}} 
= \hat{\bm{e}}\cdot \bar{\mathbf{C}} \hat{\bm{e}}.
\end{equation*}

For this to hold at all $x$, we require separate conditions at each power:
\begin{align*}
x^0: \quad & -\mathbf{X}^\top \IesanStiff \mathbf{L}_0^{-1} = \bar{\mathbf{C}}, 
\\ %\label{eq:x0-condition} \\
x^1: \quad & -\mathbf{X}^\top \mathbf{R}_1 \mathbf{L}_0^{-1} = \bar{\mathbf{C}} \mathbf{L}_0^{-1} \mathbf{L}_1 
+ \text{(symmetric partner)}.
%\label{eq:x1-condition}
\end{align*}

For the constant term \(x^0\): %\eqref{eq:x0-condition}:
\begin{equation}
\mathbf{X}^{-1} = -\bar{\mathbf{C}}^{-1} \mathbf{L}_0^{-\top} \IesanStiff^\top.
\label{eq:X-inv-formula}
\end{equation}

For the linear term \(x^1\) %\eqref{eq:x1-condition} 
to be consistent with a constant $\mathbf{X}$, 
we require the \emph{consistency condition}:
\begin{equation}
\IesanStiff^{-1} \mathbf{R}_1 = \mathbf{L}_0^{-1} \mathbf{L}_1.
\label{eq:consistency-condition}
\end{equation}

\begin{Details}

\textbf{Verification of consistency.}

The right-hand side is:
\begin{equation}
\mathbf{L}_0^{-1}\mathbf{L}_1 = \begin{pmatrix}
0 & \mathbf{0}^\top & 0 & -\CentroidLocation^\top \\
\mathbf{0} & \mathbf{0} & \mathbf{0} & \mathbf{0} \\
0 & \mathbf{0}^\top & 0 & \mathbf{0}^\top \\
\mathbf{0} & \mathbf{0} & \mathbf{0} & \oneS
\end{pmatrix}.
\end{equation}

For the left-hand side, $\mathbf{R}_1$ has only the $(4,4)$ block nonzero, 
so $\IesanStiff^{-1}\mathbf{R}_1$ acts only on the $\bm{b}_{\Section}$ columns. 
A moment increment $\delta\bm{M} = -E\FrameBasis^\times\IntRGoRG\,\bm{b}_{\Section}$ 
with $\delta N = \delta\bm{F} = \delta T = 0$ corresponds to Ieşan parameter 
increments satisfying:
\begin{align*}
0 &= EA\,\delta a_\varepsilon + EA\CentroidLocation\cdot \delta\bm{a}_\Omega, \\
\delta\bm{M} &= -EA(\FrameBasis \times \CentroidLocation)\,\delta a_\varepsilon 
- E\FrameBasis\times\IntRoR\,\delta\bm{a}_\Omega.
\end{align*}

From the first equation, $\delta a_\varepsilon = -\CentroidLocation\cdot \delta\bm{a}_\Omega$. 
Substituting into the second:
\begin{equation*}
\delta\bm{M} = -E\FrameBasis\times(\IntRoR - A\CentroidLocation \otimes \CentroidLocation)\,\delta\bm{a}_\Omega 
= -E\FrameBasis^\times\IntRGoRG\,\delta\bm{a}_\Omega.
\end{equation*}

Thus $\delta\bm{a}_\Omega = \bm{b}_{\Section}$ and 
$\delta a_\varepsilon = -\CentroidLocation\cdot \bm{b}_{\Section}$, confirming:
\begin{equation*}
\IesanStiff^{-1}\mathbf{R}_1 = \begin{pmatrix}
0 & \mathbf{0}^\top & 0 & -\CentroidLocation^\top \\
\mathbf{0} & \mathbf{0} & \mathbf{0} & \mathbf{0} \\
0 & \mathbf{0}^\top & 0 & \mathbf{0}^\top \\
\mathbf{0} & \mathbf{0} & \mathbf{0} & \oneS
\end{pmatrix} = \mathbf{L}_0^{-1}\mathbf{L}_1.
\end{equation*}
\end{Details}

\subsubsection{Trace Strain Energy}
The three-dimensional strain at section $x$ is determined by the local 
parameters $\hat{\bm{e}} = (\epsilon_0, \bm{\beta}, \hat{\varkappa}, \bm{\epsilon}_1)^\top$:
\begin{align}
\varepsilon(\bm{r}) &= \epsilon_0 + \bm{r} \cdot \bm{\epsilon}_1, \\
\bm{\varepsilon}_\gamma(\bm{r}) &= \hat{\varkappa}\,\mathbf{b}_{(\varphi)}(\bm{r}) + \mathbf{B}_{(b)}(\bm{r})\,\bm{\beta},
\end{align}
where:
\begin{equation}
\mathbf{b}_{(\varphi)} := \FrameBasis \times \bm{r} + \nabla\WarpTwistIesan, \qquad 
\mathbf{B}_{(b)} := \nabla\bm{\WarpShearIesan} + \bm{W}_{(b)}.
\end{equation}

The local parameters are related to the Ieşan parameters 
$\bm{p} = (a_\varepsilon, \bm{a}_\Omega, a_\vartheta, \bm{b}_{\Section})^\top$ by:
\begin{equation}
\hat{\bm{e}}(x) = \mathbf{L}(x)\,\bm{p} = (\mathbf{L}_0 + x\,\mathbf{L}_1)\,\bm{p},
\end{equation}
with $\mathbf{L}_0$ and \(\mathbf{L}_1\) given by \eqref{eq:L-matrices}. 
%
The section strain energy is quadratic in the local parameters:
\begin{equation}
\mathcal{U} = \frac{1}{2}\int_\Section \bm{\varepsilon}^\top \mathbb{C} \bm{\varepsilon} \, \mathrm{d}A 
= \frac{1}{2}\hat{\bm{e}} \cdot \bar{\mathbf{C}} \hat{\bm{e}},
\end{equation}
where the energy matrix takes the block form:
\begin{equation}
\bar{\mathbf{C}} = \begin{pmatrix}
EA & \mathbf{0}^\top & 0 & EA\CentroidLocation^\top \\
\mathbf{0} & G\mathbf{H}_{\psi\psi} & G\mathbf{h}_{\varphi\psi} & \mathbf{0} \\
0 & G\mathbf{h}_{\varphi\psi}^\top & \TwistRigidity & \mathbf{0}^\top \\
EA\CentroidLocation & \mathbf{0} & \mathbf{0} & E\IntRoR
\end{pmatrix},
\label{eq:energy-matrix}
\end{equation}
and
\begin{equation}
\bar{\mathbf{C}}^{-1} = \begin{pmatrix}
\frac{1}{EA} + \CentroidLocation^\top(E\IntRGoRG)^{-1}\CentroidLocation & \mathbf{0}^\top & 0 & -\CentroidLocation^\top(E\IntRGoRG)^{-1} \\[4pt]
\mathbf{0} & \frac{1}{G}\big(\mathbf{H}_{\psi\psi}^{-1} + \frac{1}{\tilde{\Jo}}\bm{\eta}\otimes\bm{\eta}\big) & -\frac{1}{G\tilde{\Jo}}\bm{\eta} & \mathbf{0} \\[4pt]
0 & -\frac{1}{G\tilde{\Jo}}\bm{\eta}^\top & \frac{1}{G\tilde{\Jo}} & \mathbf{0}^\top \\[4pt]
-(E\IntRGoRG)^{-1}\CentroidLocation & \mathbf{0} & \mathbf{0} & (E\IntRGoRG)^{-1}
\end{pmatrix}
\end{equation}

The shear-torsion coupling terms are:
\begin{align}
\mathbf{H}_{\psi\psi} &:= \int_\Section \mathbf{B}_{(b)}^\top \mathbf{B}_{(b)} \, \mathrm{d}A, \label{eq:def-H-psi-psi} \\
\mathbf{h}_{\varphi\psi} &:= \int_\Section \mathbf{B}_{(b)}^\top \mathbf{b}_{(\varphi)} \, \mathrm{d}A 
= -\Jo\CenterShearLocation + \tfrac{\nu}{2}\bm{\zeta} - \nu\bm{\lambda},
\label{eq:h-coupling}
\end{align}
where:
\begin{equation}
\bm{\zeta} := \int_\Section r^2 \nabla\WarpTwistIesan \, \mathrm{d}A, \qquad 
\bm{\lambda} := \int_\Section (\bm{r} - \CentroidLocation)(\bm{r} \cdot \nabla\WarpTwistIesan) \dA
\label{eq:def-zeta-lambda}
\end{equation}

The derivation of \eqref{eq:h-coupling} uses the boundary conditions for the 
torsion and shear warping functions: the contributions from 
$\int \nabla\bm{\WarpShearIesan} \cdot \mathbf{b}_{(\varphi)} \dA$ vanish by 
Green's identity, leaving only the Poisson warping terms.

\subsubsection{Iesan Strain Energy}
The classical resultants are related to the Ieşan parameters by 
$\bm{s} = -\mathbf{R}(x)\,\bm{p}$, where $\mathbf{R}(x) = \IesanStiff + x\mathbf{R}_1$. 
The constant part encodes the Saint-Venant constitutive relations:
\begin{equation}
\IesanStiff = \begin{pmatrix}
\AxialRigidity & EA\CentroidLocation^\top & 0 & \mathbf{0}^\top \\
\mathbf{0} & \mathbf{0} & \mathbf{0} & E\IntRGoRG \\
0 & \mathbf{0}^\top & \TwistRigidity & \mathbf{0}^\top \\
EA(\FrameBasis \times \CentroidLocation) & E\FrameBasis^\times\IntRoR & \mathbf{0} & \mathbf{0}
\end{pmatrix}
\qquad\text{and}\qquad 
\mathbf{R}_1 = \begin{pmatrix}
0 & \mathbf{0}^\top & 0 & \mathbf{0}^\top \\
\mathbf{0} & \mathbf{0} & \mathbf{0} & \mathbf{0} \\
0 & \mathbf{0}^\top & 0 & \mathbf{0}^\top \\
\mathbf{0} & \mathbf{0} & \mathbf{0} & E\FrameBasis^\times\IntRGoRG
\end{pmatrix}.
\label{eq:R0-matrix}
\end{equation}
%
where the $x$-dependent matrix \(\mathbf{R}_1\) arises from moment equilibrium, $\bm{M}' = \FrameBasis \times \bm{F}$.

\subsubsection{Strain-Amplitude Matrix.}

Computing \eqref{eq:X-inv-formula} requires inverting the block-structured 
matrices. Define the reduced torsion constant and coupling vector:
\begin{equation}
\tilde{\Jo} := \Jo - \mathbf{h}_{\varphi\psi}\cdot \mathbf{H}_{\psi\psi}^{-1} \mathbf{h}_{\varphi\psi}, 
\qquad 
\bm{\eta} := \mathbf{H}_{\psi\psi}^{-1} \mathbf{h}_{\varphi\psi}.
\label{eq:reduced-torsion}
\end{equation}

After carrying out the block matrix multiplications:
\begin{equation}
\mathbf{X}^{-1} = \begin{pmatrix}
1 & \mathbf{0}^\top & 0 & \mathbf{0}^\top \\[6pt]
\mathbf{0} & \displaystyle\frac{E}{G}\Big(\mathbf{H}_{\psi\psi}^{-1} 
+ \frac{\bm{\eta}\otimes\bm{\eta}}{\tilde{\Jo}}\Big)\IntRGoRG 
& -\Jo\mathbf{H}_{\psi\psi}^{-1}\CenterShearLocation 
- \displaystyle\frac{\Jo}{\tilde{\Jo}}\bm{\eta}(1 + \bm{\eta}\cdot\CenterShearLocation) 
& \mathbf{0} \\[6pt]
0 & -\displaystyle\frac{E}{G\tilde{\Jo}}\bm{\eta}^\top\IntRGoRG 
& \displaystyle\frac{\Jo}{\tilde{\Jo}}(1 + \bm{\eta}\cdot\CenterShearLocation) 
& \mathbf{0}^\top \\[6pt]
\mathbf{0} & \mathbf{0} & \mathbf{0} & \oneS
\end{pmatrix}.
\label{eq:X-inv-explicit}
\end{equation}


\subsubsection{Trace Matrices}

Define $(\tilde{\mathbf{A}}, \tilde{\mathbf{b}}, \tilde{\mathbf{c}}, \tilde{d})$ as the entries of the inverse of the central $3 \times 3$ block:
$$
\begin{pmatrix} \tilde{\mathbf{A}} & \tilde{\mathbf{b}} \\ \tilde{\mathbf{c}}^\top & \tilde{d} \end{pmatrix} = \begin{pmatrix} \mathbf{A} & \mathbf{b} \\ \mathbf{c}^\top & d \end{pmatrix}^{-1}
$$
where $\mathbf{A}, \mathbf{b}, \mathbf{c}, d$ are the $(2,2), (2,3), (3,2), (3,3)$ blocks of $\mathbf{X}^{-1}$. 
Then:

\begin{equation}
\boxed{\TraceDeDpRoot = \begin{pmatrix}
1 & \mathbf{0}^\top & 0 & \mathbf{0}^\top \\[4pt]
\mathbf{0} & \mathbf{0} & \tilde{\mathbf{b}} & \tilde{\mathbf{A}} + \tilde{\mathbf{b}}\CenterShearLocation^\top \\[4pt]
0 & \mathbf{0}^\top & \tilde{d} & \tilde{\mathbf{c}}^\top + \tilde{d}\CenterShearLocation^\top \\[4pt]
\mathbf{0} & \oneS & \mathbf{0} & \mathbf{0}
\end{pmatrix}
\qquad 
\TraceDeDpOne = \begin{pmatrix}
0 & \mathbf{0}^\top & 0 & -\CentroidLocation^\top \\[4pt]
\mathbf{0} & \mathbf{0} & \mathbf{0} & \mathbf{0} \\[4pt]
0 & \mathbf{0}^\top & 0 & \mathbf{0}^\top \\[4pt]
\mathbf{0} & \mathbf{0} & \mathbf{0} & \oneS
\end{pmatrix}}
\end{equation}


\subsubsection{Kinematic Reconstruction}

A key feature of the energy-conjugate trace is the relationship between 
beam-theoretic boundary conditions and the physical displacement field. 
We now examine what constraints on $\hat{\bm{u}}(x_0, \bm{r})$ are implied 
by the beam-level conditions $\bm{u}(x_0) = \bm{\theta}(x_0) = \mathbf{0}$.

Given the beam strains $\bm{e}(x) = (\bm{\gamma}(x), \bm{\kappa}(x))$, the 
kinematic fields are recovered by integration:
\begin{align}
\bm{\theta}(x) &= \bm{\theta}(0) + \int_0^x \bm{\kappa}(s)\,ds, \\
\bm{u}(x) &= \bm{u}(0) + \int_0^x \big(\bm{\gamma}(s) - \FrameBasis \times \bm{\theta}(s)\big)\,ds.
\end{align}
The six constants of integration $(\bm{u}(0), \bm{\theta}(0))$ are not 
arbitrary---they are determined by the trace applied to the underlying 
three-dimensional field.

\paragraph{The standard Ieşan solution.}

The Ieşan solution \eqref{eq:iesan-general-solution} decomposes as:
\begin{equation}
\hat{\bm{u}}(x, \bm{r}) = \underbrace{\bm{u}(x) + \bm{\theta}(x) \times \bm{r}}_{\text{section-rigid}} 
+ \underbrace{\bm{\Phi}(\bm{r})\,\bm{\alpha}(x)}_{\text{warping}},
\end{equation}
where the section-rigid part consists of polynomials in $x$ that vanish at 
$x = 0$. For the energy-conjugate trace with trace center at the origin, 
the standard Ieşan parameterization yields:
\begin{equation}
\bm{u}(0) = \mathbf{0}, \qquad \bm{\theta}(0) = \mathbf{0}.
\end{equation}

\paragraph{Physical displacement at the boundary.}

Despite $\bm{u}(0) = \bm{\theta}(0) = \mathbf{0}$, the physical displacement 
at $x = 0$ is generally nonzero:
\begin{equation}
\hat{\bm{u}}(0, \bm{r}) = \bm{\Phi}(\bm{r})\,\bm{\alpha}(0).
\end{equation}

The warping amplitude $\bm{\alpha}(0)$ includes contributions from all 
warping modes:
\begin{equation}
\hat{\bm{u}}(0, \bm{r}) = \WarpTwistIesan(\bm{r})\,\tau\,\FrameBasis 
+ \bm{W}_{(a)}(\bm{r})\,\bm{a}_\Omega 
+ \big(\bm{W}_{(b)}(\bm{r}) + \bm{\WarpShearIesan}(\bm{r})\big)\,\bm{b}_{\Section}
- \nu\bm{r}\,a_\varepsilon.
\end{equation}

For a physical point $\bm{r}^*$ to be fixed, we would require 
$\hat{\bm{u}}(0, \bm{r}^*) = \mathbf{0}$, which constitutes three scalar 
equations in two unknowns $(r_y^*, r_z^*)$. Generically, no solution exists.

\paragraph{Interpretation.}

The beam-theoretic boundary condition $\bm{u}(0) = \bm{\theta}(0) = \mathbf{0}$ 
under the energy-conjugate trace asserts:
\begin{equation}
\boxed{\text{No point of the section is fixed.}}
\end{equation}

Rather, the condition constrains only the section-rigid component of the 
motion, leaving warping entirely free. This is consistent with the 
Saint-Venant hypothesis, which precludes true Dirichlet conditions that 
would constrain warping at the boundary.

For doubly-symmetric sections with origin at the centroid, the warping 
functions satisfy $\bm{\Phi}(\mathbf{0}) = \mathbf{0}$ by symmetry, and the 
distinction disappears: $\bm{u}(0) = \mathbf{0}$ implies $\hat{\bm{u}}(0, \mathbf{0}) = \mathbf{0}$.
For asymmetric sections, however, the energy-conjugate ``clamped'' condition 
represents a fundamentally different constraint than a physical clamp.

\begin{remark}[Warping-free boundaries]
A physical clamp that fixes a region of the section would suppress warping, 
introducing boundary layer effects outside the Saint-Venant class. The 
energy-conjugate trace is incompatible with such constraints by construction. 
This is not a deficiency but rather an honest representation of what 
Saint-Venant solutions can and cannot describe.
\end{remark}

\subsubsection{Summary}

\paragraph{Properties of the energy-conjugate trace.}

\begin{description}
\item[Implementability.] 
By construction, $\bm{e} = \mathbf{X}\hat{\bm{e}}$ with $\mathbf{X}$ independent of $x$. 
The trace matrices satisfy $\TraceDeDpRoot = \mathbf{X}\mathbf{L}_0$ and 
$\TraceDeDpOne = \mathbf{X}\mathbf{L}_1$, so the implementability condition 
$\TraceDeDpOne = \TraceDeDpRoot\mathbf{L}_0^{-1}\mathbf{L}_1$ holds identically.

\item[Exactness.] 
The consistency condition \eqref{eq:consistency-condition} ensures that 
energy conjugacy holds at all sections $x$ simultaneously. The beam strains 
$\bm{e}(x)$ exactly represent the Saint-Venant solution through the local 
parameter relation $\hat{\bm{e}} = \mathbf{X}^{-1}\bm{e}$.

\item[Plane-section property.] 
The $(4,4)$ block of $\mathbf{X}^{-1}$ is $-\FrameBasis^\times$, so 
$[\mathbf{X}]_{(4,4)} = -\FrameBasis^\times$. Since 
$[\mathbf{L}_0]_{(4,2)} = \oneS$, the curvature-bending relation is:
\begin{equation}
[\TraceDeDpRoot]_{(4,2)} = [\mathbf{X}]_{(4,4)} \cdot \oneS = -\FrameBasis^\times,
\end{equation}
which matches the plane-section requirement 
$\bm{\kappa}_\perp = -\FrameBasis \times \bm{a}_\Omega$.

\item[Symmetric stiffness.] 
By the variational construction, the section stiffness matrix 
$\mathbf{C} = \mathbf{X}^\top \bar{\mathbf{C}} \mathbf{X}^{-\top}$ relating 
resultants to strains is symmetric.
\end{description}

The energy-conjugate trace thus achieves all three desiderata---implementability, 
plane-section property, and symmetric stiffness---simultaneously. This is 
possible because the consistency condition \eqref{eq:consistency-condition} 
is satisfied by the Saint-Venant structure, not imposed as an additional 
constraint.

\begin{table}[ht]
\centering
\caption{Symbols introduced in the energy-conjugate trace derivation.}
\label{tab:energy-trace-symbols}
\begin{tabular}{@{}lll@{}}
\toprule
Symbol & Units & Description \\
\midrule
\multicolumn{3}{@{}l}{\textit{Shear strain shape functions}} \\[2pt]
$\bm{\tau}_\varphi$ & L & Torsion shear strain shape, 
    $\FrameBasis \times \bm{r} + \nabla\WarpTwistIesan$ \\
$\bm{\tau}_\psi$ & L & Flexural shear strain shape, 
    $\nabla\bm{\WarpShearIesan} + \bm{W}_{(b)}$ \\[6pt]
\multicolumn{3}{@{}l}{\textit{Section stiffness quantities}} \\[2pt]
$\mathbf{H}_{\psi\psi}$ & $\mathrm{L}^4$ & Flexural shear stiffness, \eqref{eq:def-H-psi-psi} \\
$\mathbf{h}_{\varphi\psi}$ & $\mathrm{L}^4$ & Shear-torsion coupling,  \eqref{eq:reduced-torsion} \\
$\tilde{\Jo}$ & $\mathrm{L}^4$ & Reduced torsion constant, \eqref{eq:reduced-torsion} \\
$\bm{\eta}$ & --- & Shear-torsion coupling ratio, \eqref{eq:reduced-torsion} \\[6pt]
\multicolumn{3}{@{}l}{\textit{Poisson-torsion coupling} \eqref{eq:def-zeta-lambda}} \\[2pt]
$\bm{\zeta}$ & $\mathrm{L}^5$ & Torsion gradient moment, 
    $\int_\Section r^2 \nabla\WarpTwistIesan \, \mathrm{d}A$ \\
$\bm{\lambda}$ & $\mathrm{L}^5$ & Torsion-centroid coupling, 
    $\int_\Section (\bm{r} - \CentroidLocation)(\bm{r} \cdot \nabla\WarpTwistIesan) \, \mathrm{d}A$ \\[6pt]
\multicolumn{3}{@{}l}{\textit{Strain-parameter conversion}} \\[2pt]
$\mathbf{X}^{-1}$ & mixed & Beam strain to local parameter map, \eqref{eq:X-inv-explicit} \\
$\tilde{\mathbf{A}}$ & $\mathrm{L}^{-2}$ & Shear strain--flexure block of $\mathbf{X}$ \\
$\tilde{\mathbf{b}}$ & --- & Shear strain--torsion block of $\mathbf{X}$ \\
$\tilde{\mathbf{c}}$ & $\mathrm{L}^{-2}$ & Twist--flexure block of $\mathbf{X}$ \\
$\tilde{d}$ & --- & Twist--torsion block of $\mathbf{X}$ \\
\bottomrule
\end{tabular}
\end{table}

Recall: -
\(\mathbf{H} = \BishearTensorPB + \BishearTensorBB = \int \mathbf{B}_{(b)}^\top G \mathbf{B}_{(b)} \dA\)
and \(\BishearTensorRB = G\int \mathbf{B}_{(b)} \dA\) 

The correspondence is: \[\boxed{\mathbf{H} = G\,\mathbf{H}_{\psi\psi}}\]

\begin{longtable}[]{@{}
  >{\raggedright\arraybackslash}p{(\linewidth - 2\tabcolsep) * \real{0.5357}}
  >{\raggedright\arraybackslash}p{(\linewidth - 2\tabcolsep) * \real{0.4643}}@{}}
\toprule\noalign{}
\begin{minipage}[b]{\linewidth}\raggedright
Your notation
\end{minipage} & \begin{minipage}[b]{\linewidth}\raggedright
My notation
\end{minipage} \\
\midrule\noalign{}
\endhead
\bottomrule\noalign{}
\endlastfoot
\(\mathbf{H}\) & \(G\mathbf{H}_{\psi\psi}\) \\
\(\BishearTensorRB\) & \(E\IntRGoRG\) \\
\(\BfromBeta = \mathbf{H}^{-\top}\BishearTensorRB^\top\) &
\([\mathbf{X}]_{(2,2)}^{-1}\) (without torsion coupling) \\
\(\bm{K} = \frac{1}{AG}\BishearTensorRB\mathbf{H}^{-1}\BishearTensorRB^\top\)
& \(\frac{E^2}{AG^2}\IntRGoRG\mathbf{H}_{\psi\psi}^{-1}\IntRGoRG\) \\
\end{longtable}

